% Lecture 12 - Page 4 (3-Jan-1995).
% Fragment for % Lecture 12 - Page 4 (3-Jan-1995).
% Fragment for % Lecture 12 - Page 4 (3-Jan-1995).
% Fragment for % Lecture 12 - Page 4 (3-Jan-1995).
% Fragment for \input{courses/pages/curs12-pagina4-en}.
% Faithful translation in manuscript order; no preamble or forced page break.

\subsection*{Lecture 12 - Page 4}

\noindent\underline{Proof.}\quad Let $u\in A^+$,
$u=x_1x_2\cdots x_m$, $x_i\in A$.

\[
f(u)\ \overset{?}{=}\ \zeta(x_1)\zeta(x_2)\cdots\zeta(x_m)\in S.
\]
% The question mark above the equality is visible in the manuscript,
% before the proposed map is checked to be a homomorphism.

\noindent $v\in A^+$, $v=y_1y_2\cdots y_n$.
\[
f(v)=\zeta(y_1)\cdots\zeta(y_n).
\]
\[
f(uv)=\zeta(x_1)\cdots\zeta(x_m)\zeta(y_1)\cdots\zeta(y_n)
=f(u)f(v).
\]

\[
\begin{tikzcd}[column sep=large,row sep=large]
A \arrow[r,"\zeta_A"] \arrow[dr,"\zeta"'] & A^+ \arrow[d,"f"] \\
& S
\end{tikzcd}
\]

\medskip
\noindent $(!)$ Let $g:A^+\longrightarrow S$ be a semigroup homomorphism
such that $g\circ\zeta_A=\zeta$. Then $g=f$.

\medskip
\noindent $u\in A^+$,
\[
u=x_1\cdots x_m=\zeta_A(x_1)\cdots\zeta_A(x_m).
\]
\begin{align*}
g(u)
 &=g\bigl(\zeta_A(x_1)\cdots\zeta_A(x_m)\bigr)\\
 &=g\bigl(\zeta_A(x_1)\bigr)\cdots g\bigl(\zeta_A(x_m)\bigr)\\
 &=\zeta(x_1)\cdots\zeta(x_m)=f(u),
 \qquad (\forall)\,u\in A^+.
\end{align*}

\bigskip
\noindent\fbox{T.}\quad\underline{For the free monoid}

\medskip
\noindent $A\neq\varnothing$, $A^+$,
\quad $A\xrightarrow{\ \zeta_A\ }A^+$.

\medskip
\noindent For every monoid $(M,\cdot,e)$ and every set map
$\zeta:A\longrightarrow M$,
\[
\begin{tikzcd}[column sep=large,row sep=large]
A \arrow[r,"\zeta_A"] \arrow[d,"\zeta"'] & A^+ \arrow[dl,dashed,"\exists!\,f"] \\
M &
\end{tikzcd}
\]
\noindent $\exists!\,f:A^+\longrightarrow M$ such that $f\circ\zeta_A=\zeta$.

\medskip
\noindent\underline{Proof.}\quad Analogously $+$
$f(1)\overset{\mathrm{def}}{=}e\in M$.
% EDITORIAL CHECK: The heading says "free monoid" and f(1) is defined,
% but the statement and diagram use A^+ rather than A^*. The manuscript
% notation has been preserved; the free-monoid universal property would use A^*.

\bigskip
\noindent\underline{Congruences on $A^+$ generated by a binary relation}
$R\subseteq A^+\times A^+$.

\[
R=\bigl\{(w_i,w_i')\in A^+\times A^+\mid i\in I\bigr\}.
\]
% The index set appears to be I (or a similar handwritten symbol).

\noindent $(1)$\quad
$R^{-1}=\bigl\{(v,u)\mid (u,v)\in R\bigr\}$.
% The page ends here; subsequent definitions belong to the next manuscript page.
.
% Faithful translation in manuscript order; no preamble or forced page break.

\subsection*{Lecture 12 - Page 4}

\noindent\underline{Proof.}\quad Let $u\in A^+$,
$u=x_1x_2\cdots x_m$, $x_i\in A$.

\[
f(u)\ \overset{?}{=}\ \zeta(x_1)\zeta(x_2)\cdots\zeta(x_m)\in S.
\]
% The question mark above the equality is visible in the manuscript,
% before the proposed map is checked to be a homomorphism.

\noindent $v\in A^+$, $v=y_1y_2\cdots y_n$.
\[
f(v)=\zeta(y_1)\cdots\zeta(y_n).
\]
\[
f(uv)=\zeta(x_1)\cdots\zeta(x_m)\zeta(y_1)\cdots\zeta(y_n)
=f(u)f(v).
\]

\[
\begin{tikzcd}[column sep=large,row sep=large]
A \arrow[r,"\zeta_A"] \arrow[dr,"\zeta"'] & A^+ \arrow[d,"f"] \\
& S
\end{tikzcd}
\]

\medskip
\noindent $(!)$ Let $g:A^+\longrightarrow S$ be a semigroup homomorphism
such that $g\circ\zeta_A=\zeta$. Then $g=f$.

\medskip
\noindent $u\in A^+$,
\[
u=x_1\cdots x_m=\zeta_A(x_1)\cdots\zeta_A(x_m).
\]
\begin{align*}
g(u)
 &=g\bigl(\zeta_A(x_1)\cdots\zeta_A(x_m)\bigr)\\
 &=g\bigl(\zeta_A(x_1)\bigr)\cdots g\bigl(\zeta_A(x_m)\bigr)\\
 &=\zeta(x_1)\cdots\zeta(x_m)=f(u),
 \qquad (\forall)\,u\in A^+.
\end{align*}

\bigskip
\noindent\fbox{T.}\quad\underline{For the free monoid}

\medskip
\noindent $A\neq\varnothing$, $A^+$,
\quad $A\xrightarrow{\ \zeta_A\ }A^+$.

\medskip
\noindent For every monoid $(M,\cdot,e)$ and every set map
$\zeta:A\longrightarrow M$,
\[
\begin{tikzcd}[column sep=large,row sep=large]
A \arrow[r,"\zeta_A"] \arrow[d,"\zeta"'] & A^+ \arrow[dl,dashed,"\exists!\,f"] \\
M &
\end{tikzcd}
\]
\noindent $\exists!\,f:A^+\longrightarrow M$ such that $f\circ\zeta_A=\zeta$.

\medskip
\noindent\underline{Proof.}\quad Analogously $+$
$f(1)\overset{\mathrm{def}}{=}e\in M$.
% EDITORIAL CHECK: The heading says "free monoid" and f(1) is defined,
% but the statement and diagram use A^+ rather than A^*. The manuscript
% notation has been preserved; the free-monoid universal property would use A^*.

\bigskip
\noindent\underline{Congruences on $A^+$ generated by a binary relation}
$R\subseteq A^+\times A^+$.

\[
R=\bigl\{(w_i,w_i')\in A^+\times A^+\mid i\in I\bigr\}.
\]
% The index set appears to be I (or a similar handwritten symbol).

\noindent $(1)$\quad
$R^{-1}=\bigl\{(v,u)\mid (u,v)\in R\bigr\}$.
% The page ends here; subsequent definitions belong to the next manuscript page.
.
% Faithful translation in manuscript order; no preamble or forced page break.

\subsection*{Lecture 12 - Page 4}

\noindent\underline{Proof.}\quad Let $u\in A^+$,
$u=x_1x_2\cdots x_m$, $x_i\in A$.

\[
f(u)\ \overset{?}{=}\ \zeta(x_1)\zeta(x_2)\cdots\zeta(x_m)\in S.
\]
% The question mark above the equality is visible in the manuscript,
% before the proposed map is checked to be a homomorphism.

\noindent $v\in A^+$, $v=y_1y_2\cdots y_n$.
\[
f(v)=\zeta(y_1)\cdots\zeta(y_n).
\]
\[
f(uv)=\zeta(x_1)\cdots\zeta(x_m)\zeta(y_1)\cdots\zeta(y_n)
=f(u)f(v).
\]

\[
\begin{tikzcd}[column sep=large,row sep=large]
A \arrow[r,"\zeta_A"] \arrow[dr,"\zeta"'] & A^+ \arrow[d,"f"] \\
& S
\end{tikzcd}
\]

\medskip
\noindent $(!)$ Let $g:A^+\longrightarrow S$ be a semigroup homomorphism
such that $g\circ\zeta_A=\zeta$. Then $g=f$.

\medskip
\noindent $u\in A^+$,
\[
u=x_1\cdots x_m=\zeta_A(x_1)\cdots\zeta_A(x_m).
\]
\begin{align*}
g(u)
 &=g\bigl(\zeta_A(x_1)\cdots\zeta_A(x_m)\bigr)\\
 &=g\bigl(\zeta_A(x_1)\bigr)\cdots g\bigl(\zeta_A(x_m)\bigr)\\
 &=\zeta(x_1)\cdots\zeta(x_m)=f(u),
 \qquad (\forall)\,u\in A^+.
\end{align*}

\bigskip
\noindent\fbox{T.}\quad\underline{For the free monoid}

\medskip
\noindent $A\neq\varnothing$, $A^+$,
\quad $A\xrightarrow{\ \zeta_A\ }A^+$.

\medskip
\noindent For every monoid $(M,\cdot,e)$ and every set map
$\zeta:A\longrightarrow M$,
\[
\begin{tikzcd}[column sep=large,row sep=large]
A \arrow[r,"\zeta_A"] \arrow[d,"\zeta"'] & A^+ \arrow[dl,dashed,"\exists!\,f"] \\
M &
\end{tikzcd}
\]
\noindent $\exists!\,f:A^+\longrightarrow M$ such that $f\circ\zeta_A=\zeta$.

\medskip
\noindent\underline{Proof.}\quad Analogously $+$
$f(1)\overset{\mathrm{def}}{=}e\in M$.
% EDITORIAL CHECK: The heading says "free monoid" and f(1) is defined,
% but the statement and diagram use A^+ rather than A^*. The manuscript
% notation has been preserved; the free-monoid universal property would use A^*.

\bigskip
\noindent\underline{Congruences on $A^+$ generated by a binary relation}
$R\subseteq A^+\times A^+$.

\[
R=\bigl\{(w_i,w_i')\in A^+\times A^+\mid i\in I\bigr\}.
\]
% The index set appears to be I (or a similar handwritten symbol).

\noindent $(1)$\quad
$R^{-1}=\bigl\{(v,u)\mid (u,v)\in R\bigr\}$.
% The page ends here; subsequent definitions belong to the next manuscript page.
.
% Faithful translation in manuscript order; no preamble or forced page break.

\subsection*{Lecture 12 - Page 4}

\noindent\underline{Proof.}\quad Let $u\in A^+$,
$u=x_1x_2\cdots x_m$, $x_i\in A$.

\[
f(u)\ \overset{?}{=}\ \zeta(x_1)\zeta(x_2)\cdots\zeta(x_m)\in S.
\]
% The question mark above the equality is visible in the manuscript,
% before the proposed map is checked to be a homomorphism.

\noindent $v\in A^+$, $v=y_1y_2\cdots y_n$.
\[
f(v)=\zeta(y_1)\cdots\zeta(y_n).
\]
\[
f(uv)=\zeta(x_1)\cdots\zeta(x_m)\zeta(y_1)\cdots\zeta(y_n)
=f(u)f(v).
\]

\[
\begin{tikzcd}[column sep=large,row sep=large]
A \arrow[r,"\zeta_A"] \arrow[dr,"\zeta"'] & A^+ \arrow[d,"f"] \\
& S
\end{tikzcd}
\]

\medskip
\noindent $(!)$ Let $g:A^+\longrightarrow S$ be a semigroup homomorphism
such that $g\circ\zeta_A=\zeta$. Then $g=f$.

\medskip
\noindent $u\in A^+$,
\[
u=x_1\cdots x_m=\zeta_A(x_1)\cdots\zeta_A(x_m).
\]
\begin{align*}
g(u)
 &=g\bigl(\zeta_A(x_1)\cdots\zeta_A(x_m)\bigr)\\
 &=g\bigl(\zeta_A(x_1)\bigr)\cdots g\bigl(\zeta_A(x_m)\bigr)\\
 &=\zeta(x_1)\cdots\zeta(x_m)=f(u),
 \qquad (\forall)\,u\in A^+.
\end{align*}

\bigskip
\noindent\fbox{T.}\quad\underline{For the free monoid}

\medskip
\noindent $A\neq\varnothing$, $A^+$,
\quad $A\xrightarrow{\ \zeta_A\ }A^+$.

\medskip
\noindent For every monoid $(M,\cdot,e)$ and every set map
$\zeta:A\longrightarrow M$,
\[
\begin{tikzcd}[column sep=large,row sep=large]
A \arrow[r,"\zeta_A"] \arrow[d,"\zeta"'] & A^+ \arrow[dl,dashed,"\exists!\,f"] \\
M &
\end{tikzcd}
\]
\noindent $\exists!\,f:A^+\longrightarrow M$ such that $f\circ\zeta_A=\zeta$.

\medskip
\noindent\underline{Proof.}\quad Analogously $+$
$f(1)\overset{\mathrm{def}}{=}e\in M$.
% EDITORIAL CHECK: The heading says "free monoid" and f(1) is defined,
% but the statement and diagram use A^+ rather than A^*. The manuscript
% notation has been preserved; the free-monoid universal property would use A^*.

\bigskip
\noindent\underline{Congruences on $A^+$ generated by a binary relation}
$R\subseteq A^+\times A^+$.

\[
R=\bigl\{(w_i,w_i')\in A^+\times A^+\mid i\in I\bigr\}.
\]
% The index set appears to be I (or a similar handwritten symbol).

\noindent $(1)$\quad
$R^{-1}=\bigl\{(v,u)\mid (u,v)\in R\bigr\}$.
% The page ends here; subsequent definitions belong to the next manuscript page.
